\documentclass[11pt]{article}
\usepackage[margin=1in]{geometry}
\usepackage{amsmath,amssymb}
\usepackage{booktabs}
\usepackage{slashed}
\usepackage[colorlinks=true,linkcolor=blue,urlcolor=blue]{hyperref}
\usepackage{parskip}
\setlength{\parindent}{0pt}

\title{\vspace{-2em}Comments on\\[0.3em]
\emph{Fully accounted NLO electromagnetic radiative corrections\\
to the width of $J/\psi$ decay}}
\author{V. Kubarovsky}
\date{16 September 2026}

\begin{document}
\maketitle

\noindent
Equation and table numbers refer to the current draft (\texttt{w\_jpsi-4}).
``VPK note'' refers to \emph{Radiative Corrections to the $J/\psi$
photoproduction reaction}.

\section{Eq.~(1) and Eq.~(17): replace $E_V$ by $m_V$}

Both read $1/(2E_V)$. The whole calculation is done in the $J/\psi$ rest frame,
where $E_V=m_V$ identically, and the text already says so. Writing $m_V$ removes
a symbol that never takes any other value, and makes the normalisation of
Eq.~(1) directly comparable with the standard two-body form.

\section{Eq.~(6): use the factorised form, and add the ultrarelativistic limit}

Eq.~(6) is correct --- it is algebraically identical to Eq.~(7) of the VPK note
(verified symbolically; the difference is exactly zero). But the factorised form
is easier to read and to check dimensionally, because the mass dependence enters
only through dimensionless ratios:
\begin{equation*}
\Gamma_B=\frac{4\pi\alpha^{2}}{f_V^{2}}\,\frac{m_V}{3}
\left(1+\frac{2m_\ell^{2}}{m_V^{2}}\right)
\sqrt{1-\frac{4m_\ell^{2}}{m_V^{2}}}\,.
\end{equation*}
Please also add the ultrarelativistic limit, Eq.~(8) of the VPK note, which is
the form actually used to fix $f_V$ from the measured leptonic width:
\begin{equation*}
\Gamma_B=\frac{4\pi\alpha^{2}}{f_V^{2}}\,\frac{m_V}{3}\,.
\end{equation*}
Worth stating explicitly: the correction connecting the two is \emph{second}
order in the mass ratio, not first. With $x=m_\ell^{2}/m_V^{2}$,
\begin{equation*}
(1+2x)\sqrt{1-4x}=1-6x^{2}+O(x^{3}),
\end{equation*}
so the linear term cancels. This is why the massless limit is so accurate:
$4\times10^{-15}$ for $J/\psi\to e^+e^-$ and $8\times10^{-6}$ for
$J/\psi\to\mu^+\mu^-$. It reaches only $2\times10^{-3}$ for $\rho\to\mu^+\mu^-$.

\section{Table~I: please say what $m_V=1.5$~GeV is}

Tables~I and~II are computed at $m_V=1.5$~GeV, which is neither the $J/\psi$ nor
any physical vector meson. A reader meeting it just after Section~II, which sets
$m_V=3.0969$~GeV, will be confused.

We understand it to be the last row of Table~I of Ref.~[10], whose tabulated
column is $2E$~(MeV), running from 500 to 1500 with $\rho$, $\omega$ and $\phi$
marked --- that is, $2E$ is the meson mass, since $E=m_V/2$ at resonance (stated
in [10] below their Eq.~(22)). One sentence would settle it: ``for comparison
with Ref.~[10] we take $m_V=1.5$~GeV, the largest mass tabulated there.''

\section{The comparison with Ref.~[10] does not hold where it is claimed}

The text states that Table~I can be ``successfully compared'' with Table~I of
Ref.~[10] at $\omega/m_e=0.01$, last line of the first half. At that line the two
disagree, and the disagreement grows with $r_{cut}$:

\begin{center}
\begin{tabular}{cccc}
\toprule
$r_{max}$ & Ref.~[10] & this work & difference\\
\midrule
0.10 & $-0.087$ & $-0.0866$    & $\sim 0$\\
0.25 & $-0.034$ & $-0.0330$    & $0.001$\\
0.50 & $-0.002$ & $+0.000044$  & $0.002$\\
1.00 & $+0.021$ & $+0.023891$  & $0.0029$\\
\bottomrule
\end{tabular}
\end{center}

The soft part agrees exactly (both give $\delta_R=-0.747$), so the mass and the
cut-off are matched correctly. The discrepancy is entirely in the hard part and
equals $5\alpha/4\pi=0.0029$ at $r_{max}=1$.

The reason is that \textbf{Eq.~(24) of Ref.~[10] is missing a term}. Their
differential spectrum, Eq.~(23), is correct and agrees with Eq.~(14) of this
work; the term was lost when that spectrum was integrated analytically. Their
Eq.~(26), and their Table~I, inherit the omission. Recomputing their table from
their own Eq.~(26) reproduces every printed digit, electrons and muons alike,
which confirms the diagnosis.

\textbf{Where to record this.} In Section~IV of the paper, immediately after the
closed-form expressions we propose in Section~5 of these comments --- i.e.\ right
where the new formulas will appear, following your Eq.~(14). Those expressions
are precisely Eq.~(26) of [10] with the term restored, so a single paragraph
placed there does three things at once: it attributes the formula to [10],
records what had to be corrected, and explains why the numerical comparison with
Table~I of [10] holds only in the small-$r_{cut}$ rows. The provenance and the
correction then appear together, where a reader using the formula will actually
meet them.

The sentence after Table~I claiming the comparison ``succeeds'' would become a
pointer to that paragraph. Suggested wording:

\begin{quote}
The expressions above are Eq.~(26) of Ref.~[10] with the term
$(\alpha/\pi)(3r_{max}/2-r_{max}^{2}/4)$ restored. The differential spectrum,
Eq.~(23) of [10], is correct and agrees with our Eq.~(14); the term was lost in
its analytic integration, so that Eq.~(24), Eq.~(26) and Table~I of [10] are all
low by that amount. The deficit is $5\alpha/4\pi=0.0029$ at $r_{max}=1$ and falls
below the precision quoted in [10] for $r_{max}\lesssim0.1$, which is why our
results reproduce the small-$r_{max}$ rows of that table but not the last one.
\end{quote}

\section{Please add the closed-form total correction, for both channels}

The paper gives $\delta_{self}$~(8), $\delta_{vert}$~(12), $\delta_{soft}$~(13),
and $\delta_{hard}$ only as an unevaluated integral~(14), then the exact
result~(28). There is no compact closed form for the total correction, which is
what an experiment actually applies and what a reader comes to this paper
looking for. We would like it in the paper, written out separately for the two
channels. These are Eqs.~(19) and~(20) of the VPK note, i.e.\ Eq.~(26) of
Ref.~[10] with the missing term restored.

\medskip
\textbf{Electrons}, $J/\psi\to e^+e^-$:
\begin{equation*}
\boxed{\;
\delta(m_V,k_{max})=-\frac{4\alpha}{\pi}
\left[
\left(\frac12-\ln\frac{m_V}{m_e}\right)\left(\ln r_{max}+\frac13\right)
+\frac12\left(L_2(r_{max})-\frac{\pi^{2}}{6}\right)
+\frac{17}{72}+R(r_{max})
+\frac{r_{max}^{2}}{16}-\frac{3}{8}r_{max}
\right]\;}
\end{equation*}

\textbf{Muons}, $J/\psi\to\mu^+\mu^-$:
\begin{equation*}
\boxed{\;
\begin{aligned}
\delta(m_V,k_{max})=-\frac{4\alpha}{\pi}
\Bigg[&
\left(\frac12-\ln\frac{m_V}{m_\mu}\right)\left(\ln r_{max}+\frac13\right)
+\frac12\left(L_2(r_{max})-\frac{\pi^{2}}{6}\right)
+\frac{17}{72}\\[2pt]
&+\left(\frac{5}{18}-\frac13\ln\frac{m_V}{m_e}\right)
+R(r_{max})
+\frac{r_{max}^{2}}{16}-\frac{3}{8}r_{max}
\Bigg]
\end{aligned}
\;}
\end{equation*}
with, in each case,
\begin{align*}
r_{max}&=\frac{2k_{max}}{m_V}, &
R(r_{max})&=\left(r_{max}-\frac{r_{max}^{2}}{4}-\frac34\right)
\left[\ln\!\left(\frac{m_V}{m_\ell}\sqrt{1-r_{max}}\right)-\frac12\right], &
L_2(x)&=-\int_0^x\frac{\ln(1-t)}{t}dt,
\end{align*}
$m_\ell$ being the mass of the lepton in the channel concerned. Here $r_{max}$
is Ehlotzky's variable, defined below his Eq.~(25) as ``the maximum fraction of
the lepton energy $E$ \ldots which is lost \ldots by the emission of a hard
photon'', i.e.\ $r_{max}=k_{max}/E$ with $E=m_V/2$. It is identical to the
$r_{cut}=\omega_{cut}/E_1$ already defined in this paper, so no new variable is
needed --- but see Section~6. The extra term
in the muon expression, $\frac{5}{18}-\frac13\ln\frac{m_V}{m_e}$, is the
electron-loop vacuum polarization; it is exactly $-\frac{\pi}{4\alpha}$ times
$\delta_{self}$ of Eq.~(8) evaluated with $m_e$, and it correctly keeps $m_e$,
not $m_\mu$.

The last two terms of each bracket, $r_{max}^{2}/16-\tfrac38 r_{max}$, are the
ones absent from Ref.~[10].

\subsection*{Where it belongs}

Section~IV, immediately after Eq.~(14) --- the closed form \emph{is} the
evaluation of that integral, combined with Eqs.~(8), (12) and~(13). Three
reasons it sits naturally there and nowhere else:

\begin{itemize}
\item Section~IV opens by promising ``compact expressions for all contributions
      to decay width: both for virtual and real bremstrahlung ones.'' Compact
      expressions are then given for every individual contribution, but not for
      the sum, which is the only combination that is physical.
\item Eq.~(16), $\delta_{FSIC}=3\alpha/4\pi$, becomes an immediate corollary: it
      is the $r_{max}\to1$ limit of the closed form, less $\delta_{self}$.
\item Figure~4 plots exactly this quantity for both channels. At present the
      reader is shown the curve without being given the function.
\end{itemize}

One caveat on placement. Section~IV states that the URA is ``applicable for
$\ell=e$'', and that the mass-dependent treatment is ``crucial for \ldots the
muon channel''. That is consistent with the numbers below: the electron formula
is exact for all practical purposes, the muon one is good to $\sim10^{-4}$. We
would therefore present the electron expression as the working formula of
Section~IV, and the muon expression alongside it with an explicit note that it
carries the $(m_\mu/E)^{2}$ error quantified below --- pointing the reader to
Eq.~(28) when that matters.

\subsection*{Status of the approximation}

These expressions are \emph{not} exact in the lepton mass, and they are also not
the $m_\ell=0$ limit. They are the ultrarelativistic result: the lepton mass is
kept in the logarithms, while terms suppressed by $(m_\ell/E)^{2}$, $E=m_V/2$,
are dropped. Ref.~[10] states this explicitly just above its Eq.~(22)
(``\ldots neglecting terms of order $(m/E)^{2}$''), and the same restriction
propagates to its Eqs.~(24) and~(26). We suggest saying so when the formulas are
introduced, since this paper's whole point is the mass-exact treatment.

The size of what is dropped:
\begin{center}
\begin{tabular}{lccc}
\toprule
 & $(m_\ell/E)^{2}$ & $\delta_{self}$: Eq.~(8) $-$ Eq.~(9) & total at $r_{max}=1$: this form $-$ Eq.~(28)\\
\midrule
$e^+e^-$     & $1.1\times10^{-7}$ & $+3\times10^{-10}$ & $+1\times10^{-7}$\\
$\mu^+\mu^-$ & $4.7\times10^{-3}$ & $+1.1\times10^{-5}$ & $-7\times10^{-5}$\\
\bottomrule
\end{tabular}
\end{center}
So for the electron channel the closed form is exact at any precision an
experiment will reach, and for the muon channel it is good to $\sim10^{-4}$
absolute --- which is precisely the accuracy argument made in Section~IV of the
draft. Presenting the closed form alongside Eq.~(28) would let a reader use the
simple expression where it suffices and the exact one where it does not.

Checks the expressions pass:
\begin{itemize}
\item they reproduce Table~I of this work to all six digits at every $r_{cut}$,
      in both channels, and are independent of $\omega$ by construction;
\item at $r_{max}=1$ they give exactly $\delta_{self}+3\alpha/4\pi$, i.e.\
      $\delta_{self}$ plus the FSIC of Eq.~(16);
\item for $J/\psi\to e^+e^-$ at $r_{max}=1$ the value is $0.026135$, matching the
      Bardin--Shumeiko result of Table~III.
\end{itemize}

\section{The upper limit of $r_{cut}$ is not 1}

The text before Eq.~(16) says the cut energy takes its maximum value ``at (15) or
$r_{cut}=r_{max}=1$'', and Table~III is captioned $(r_{max}=1)$. That is
inconsistent with Eq.~(15) itself. With $r_{cut}=\omega_{cut}/E_1=2\omega_{cut}/m_V$
and $\omega_{max}=m_V/2-2m_\ell^{2}/m_V$,
\begin{equation*}
r_{cut}^{\;\max}=\frac{2\omega_{max}}{m_V}=1-\frac{4m_\ell^{2}}{m_V^{2}}=\beta^{2},
\end{equation*}
which is $0.99535$ for muons, not 1. (For electrons it is $1-1.1\times10^{-7}$.)

This is not only notational. Eq.~(29) defines
$\rho_1(\omega_{cut})=\sqrt{m_V(m_V-2\omega_{cut})}+\sqrt{m_V(m_V-2\omega_{cut})-4m_\ell^{2}}$.
Setting $\omega_{cut}=m_V/2$, i.e.\ literally $r_{cut}=1$, makes the second
radicand $-4m_\ell^{2}<0$ and $\rho_1$ imaginary. At the true endpoint
$\omega_{cut}=\omega_{max}$ that radicand vanishes exactly and
$\rho_1(\omega_{max})=2m_\ell$, which is the correct and finite limit. So Eq.~(28)
is well defined up to $\omega_{max}$ and no further.

We take it that Table~III was in fact evaluated at $\omega_{cut}=\omega_{max}$,
which is the sensible thing to do; only the label is wrong. We suggest replacing
``$r_{cut}=r_{max}=1$'' and the Table~III caption by ``$\omega_{cut}=\omega_{max}$''
or ``$r_{cut}=\beta^{2}$''. The same remark applies to the upper end of the axis
in Fig.~4, which is correctly drawn as $\omega_{cut}/\omega_{max}$.

\section{Reference [9]}

\begin{center}
\texttt{[9] P. Chatagnon, e-Print: 2602.22128 [hep-ex] (2026).}
\end{center}
is now published. Please replace with
\begin{center}
\texttt{[9] P. Chatagnon, Phys. Rev. C \textbf{113} (2026) 6, 065203.}
\end{center}

\section*{Additional point, for consideration}

\textbf{Eq.~(4).} As printed,
\begin{equation*}
\sum_{\rm pol}\varepsilon_\alpha(p_V)\varepsilon_\beta(p_V)=-g_{\alpha\beta}
\end{equation*}
is the polarization sum for a \emph{massless} vector. For a massive vector meson
it is
\begin{equation*}
\sum_\lambda\varepsilon_\alpha\varepsilon_\beta^{*}
=-g_{\alpha\beta}+\frac{p_\alpha p_\beta}{m_V^{2}}\,.
\end{equation*}
The extra term does drop here, because it contracts the lepton current with
$p_V=p_1+p_2$ and $\bar u(p_1)(\slashed p_1+\slashed p_2)v(p_2)=0$ by the Dirac
equations --- so Eq.~(5) and everything after it are correct. But the identity as
written is not true on its own, and the surrounding text calls it ``the photon''
while the sum is over the three $J/\psi$ polarizations. A clause such as ``the
$p_\alpha p_\beta$ term is dropped by current conservation'' would fix it.

\end{document}
